\documentclass[10pt,a4paper]{article}
\usepackage[T1]{fontenc}
\usepackage[utf8]{inputenc}
\usepackage{lmodern}
\usepackage{amsmath,amssymb}
\usepackage[margin=21mm]{geometry}
\usepackage{microtype}
\usepackage[hidelinks]{hyperref}
\usepackage{enumitem}
\setlist[itemize]{leftmargin=*,nosep,topsep=3pt}
\setlength{\parindent}{0pt}
\setlength{\parskip}{5pt}
\setlength{\emergencystretch}{3em}
\newcommand{\question}[1]{\par\medskip{\normalsize\bfseries #1}\par\smallskip}
\pagestyle{plain}

\begin{document}
{\Large\bfseries Optogenetics\par}

\textbf{Paper:} Karl Deisseroth (2011).\quad Foundational perspective; author manuscript.\par
\href{https://doi.org/10.1038/nmeth.f.326}{DOI: \nolinkurl{10.1038/nmeth.f.326}}\par
{\footnotesize Page references below use PDF page numbers, including front matter.}

\question{To which general question does this paper relate?}
What optogenetics can control and how selective perturbations can test causal circuit hypotheses. Provides terminology rather than a disease-specific feedback objective.

\question{How does it answer the general questions?}
\textbf{Actuation.} Genetically targeted light-sensitive tools enable rapid excitation, inhibition and modulation of defined cells or projections. Tool choice changes temporal precision, light sensitivity and the duration of induced excitability changes (PDF pp.~1--5, 9).

\textbf{Measurement and causality.} Electrical recordings or optical readouts must separately establish the resulting neural activity. Selective interventions can test whether a cell class or projection contributes to a phenotype, whereas observing its activity alone cannot establish that contribution. The article also stresses that reliable targeting and validated opsin function matter (PDF pp.~2--6).

\textbf{Model relevance.} Supports selecting where $u$ enters a circuit and checking its sign and dynamics experimentally. It gives no unique light-to-firing gain, delay, sigmoid slope or measurement matrix for the notebook. Millisecond-scale capability should not be treated as a universal actuator time constant. No beta-suppression objective or PIE multiplier design follows from this perspective.

\question{What papers does it cite that could provide more information?}
\begin{itemize}
\item Boyden et al. (2005), \emph{Millisecond-timescale, genetically targeted optical control of neural activity}: foundational neuronal actuation (ref.~4, PDF pp.~6--7).
\item Gradinaru et al. (2009), \emph{Optical deconstruction of parkinsonian neural circuitry}: circuit-specific Parkinsonian interventions (ref.~19, PDF p.~7).
\end{itemize}

\end{document}
